\chapter{User research per la Evaluation}
\label{ch:userresearch}
La User Evaluation d'esame richiede metodo, profilo, risultati e riflessioni.
Seguiamo Bassanelli: metodi \textbf{qualitativi} (think-aloud durante boss
fight e deck building, interviste semi-strutturate) + \textbf{quantitativi}
(questionari: SUS per l'usabilità, IMI per la motivazione intrinseca,
GEQ-core per l'esperienza di gioco).

\section{Protocollo (da eseguire)}
\begin{enumerate}[nosep]
\item Reclutamento: 5--8 studenti con Flutter di base (profilo da descrivere).
\item Sessione (30\,min): 5\,min onboarding, 15\,min gioco guidato
(roadmap $\rightarrow$ quiz $\rightarrow$ boss) con think-aloud,
10\,min questionari + intervista.
\item Metriche: tasso superamento quiz, tempo/boss, carte usate, SUS/IMI/GEQ.
\item Riflessioni: cosa conferma/smentisce Octalysis (CD3/CD7) e Toda (gap).
\end{enumerate}
Nessun dato utente è stato ancora raccolto: nel GamiDOC la sezione resta un
placeholder esplicito, senza inventare risultati.
